
In this section, we provide details of the coreference systems and the corpus we analyze. We consider three systems with different types of techniques, rule-based, feature-based, and neural network-based. We describe these systems and the corpus below. 
 
\textbf{The Rule-based System.}
The Stanford Deterministic Coref System uses several layers of rules for making coreference decisions on an article. For each layer, the model read the article from the beginning to the end, and for each mention, it decides to associate the mention with some mentions appeared in the previous context or create a new cluster. The decisions are based on a set of syntactic and semantic rules that generated in an unsupervised fashion. For example, the rules include  identifying  \textit{exact match} between mentions and to match mentions based on their \textit{head word}\footnote{The word associate with the root in the syntactic parse tree of the mention. For example, the head word of ``a Ph.D. student'' is ``student''.} and their attributes such as gender, number, and animacy. 
Despite the model is simple and only based on rules, it is the winner for the CoNLL shared task 2011 on the coreference resolution task.

\textbf{The Feature-based System.}
The Berkeley coref system is a learning-based system~\cite{DurrettKlein2013} which aims at solving the coreference resolution using features extracted from the input. The model captures both linguistic features as well as underlying data patterns which are not easily predefined by rules. The model applies mention ranking approach~\cite{denis2008specialized} with using logistic regression as the base learner. The features used in the system include lexicon features such as 
pairs of head words in both mentions. Therefore, the model will automatically learn the association between the coreference and the occurrence between words in the mentions. For example, the model will learn to associate male with president when the training data contains more cases where ``he'' refers to ``the president''.
 
\textbf{The Neural-based System.}
The UW coref system is an end-to-end coreference resolution model which employs a neural network that directly learns a representation of entities and mentions from text. Different from the feature-based approaches, the model does not require hand-designed features and
syntactic parse trees. Instead, the mention boundaries and the coreference decisions can be directly learned from the data.
At the core of the model are vector embeddings to represent spans of text in  documents. These embeddings combine the context-dependent lexical information extracted by a bi-direction LSTM  with the  head-finding attention mechanism.
This end-to-end neural-based model currently achieves the best performance on the OntoNotes 5.0 corpus. 

\paragraph{\bf The OntoNotes 5.0 Corpus}
The analyses in this paper are based on the English portion of OntoNotes 5.0 corpus~\cite{pradhan2012conll}. This corpus is used in the CoNLL-2012 shared task, the main modeling competition in the natural language processing community. The corpus consists of text collected from seven different sources, including Broadcast Conversation, Broadcast News, Magazine, Newswire, Telephone Conservation, Weblogs and Newsgroup and Pivot Text.   
The corpus consists 1.6 million words in total coming from  $2,384$ documents. We use the standard data split based on the CoNLL 2012 shared task~\cite{pradhan2012conll}, and partition the corpus into 1,940, 222, 222, documents for training, development, and test, respectively. 
Note that some documents are very long; therefore, the annotators split them into several parts and only consider annotating coreference chains within each part.  
For coreference resolution tasks, the model takes documents to make predictions for event and entity mentions and to link the related mentions together to form the coreference chains. To evaluate the performance of the models,  we follow the evaluation metrics used in CoNLL-2012 shared task\footnote{\url{http://conll.cemantix.org/2012/software.html}} and use the official scorer provided~\cite{luo2014extension,pradhan2014scoring}.